% !TeX root = ../../Book3.tex
\section{Green 域与抛物型域}
\label{b3:sec:1702}

上一节在有界正则域中把点源与连续零边值配在一起。若区域无界，或者边界含有不正则点，就不能继续把通常的边界值当作唯一的约束。本节先问一个更基本的问题：在整个区域内，能否有一个始终为正、只有指定对数奇性的函数？下一节再从所有这样的函数中构造最小者。

\subsection{Green 域}
\label{b3:subsec:170201}

为了在尚未讨论无穷远行为时准确表述存在性，先把竞争函数类写出来。

\Needspace{6\baselineskip}
\begin{definition}[正基本核与 Green 域]\label{b3:def:17-green-domain}
设 \(D\subseteq\mathbb C\) 为区域，\(w\in D\)。称实函数 \(q:D\setminus\{w\}\to(0,\infty)\) 为极点 \(w\) 的正基本核，如果它在 \(D\setminus\{w\}\) 调和，且 \(q(z)+\log|z-w|\) 在 \(w\) 的某个圆盘邻域中延拓为调和函数。记所有正基本核的集合为 \(\mathcal Q_D(w)\)。若对每个 \(w\in D\) 都有 \(\mathcal Q_D(w)\ne\varnothing\)，称 \(D\) 为\term[green-yu]{Green 域}[Greenian domain]。当 \(\mathcal Q_D(w)\) 中存在逐点不超过其余各项的成员时，称这个最小成员为一般域上的 Green 核，记为 \(g_D(\cdot,w)\)。
\end{definition}

“最小”指与全体竞争函数逐点比较，不是只在某个内点使数值最小。一个类如果有最小成员，该成员必唯一。上一节已经证明，有界正则域的经典 Green 核就是其正基本核类的最小成员。下一节将证明，只要该类非空，最小成员就存在；还将证明一处极点的存在性足以推出全部极点的存在性。因此这里的定义与用有限穷竭极限定义 Green 域等价，但现在不把尚未构造的极限作为论证前提。

\ProseParagraphBreak
有界域不必正则，也一定有正基本核。固定 \(w\in D\)，选择 \(R>0\) 使 \(D\subset B(w,R)\)，则
\[
 q(z)=\log\frac R{|z-w|}
\]
在 \(D\setminus\{w\}\) 严格为正、调和，且 \(q+\log|z-w|=\log R\)。所以每个有界区域都是 Green 域。这个竞争核一般不在 \(\partial D\) 上取零；其用途是给随后构造的最小核一个有限上界。边界是否正则影响连续零边值的达到，不决定有界域是否有 Green 核。

\ProseParagraphBreak
正基本核也可以限制到子域。若 \(D\subset E\)，\(w\in D\)，则 \(\mathcal Q_E(w)\) 的每一项在 \(D\) 上仍有正确奇性和正性。因此 Green 域的每个子域也是 Green 域。这一判断只用竞争函数的存在，尚不需要最小核的区域单调性。

\subsection{Green 函数不存在的情形}
\label{b3:subsec:170202}

全平面的困难出现在无穷远。假设极点 \(w\) 有一个正基本核 \(q\)，则
\[
 h(z)=q(z)+\log|z-w|
\]
延拓为整个平面的调和函数。在以 \(w\) 为中心的圆周上取平均，得到
\[
 \frac1{2\pi}\int_0^{2\pi}q(w+re^{\mathrm it})\,\mathrm dt
 =h(w)-\log r.
\]
当 \(r\) 足够大时右端为负，与 \(q>0\) 矛盾。因此整个复平面没有正基本核，更没有 Green 核。这里不是因为没有可写的基本解：\(-\log|z-w|\) 正是所需基本解，只是任何整平面调和修正都不能使它在全部半径上保持正平均。

\ProseParagraphBreak
删除有限个点是否可以补救这一困难，需要先控制那些点附近可能出现的奇性。对于任意调和函数，孤立奇性可以比对数强得多；正性却排除了这些振荡主部。

\Needspace{6\baselineskip}
\begin{lemma}[Bôcher 孤立奇性分解]\label{b3:lem:17-positive-singularity}
设 \(R>0\)，\(u\ge0\) 在 \(0<|z|<R\) 调和。则存在唯一的 \(c\ge0\) 及整个圆盘 \(B(0,R)\) 上的实调和函数 \(h\)，使
\[
 u(z)=-c\log|z|+h(z)\quad(0<|z|<R).
\]
\end{lemma}
\begin{proof}
对每个整数 \(n\) 定义圆周系数
\[
 a_n(r)=\frac1{2\pi}\int_0^{2\pi}u(re^{\mathrm it})e^{-\mathrm int}\,\mathrm dt,
 \quad 0<r<R.
\]
任意闭圆环内的 \(u\) 光滑，故可以对 \(r\) 求导；将极坐标 Laplace 方程乘以 \(e^{-\mathrm int}\)，并对角变量积分两次分部，得到
\[
 a_n''(r)+\frac1r a_n'(r)-\frac{n^2}{r^2}a_n(r)=0.
\]
当 \(n=0\) 时，解为 \(a_0(r)=A+B\log r\)。圆周平均非负，令 \(r\downarrow0\) 可知 \(B\le0\)。当 \(n\ne0\) 时，换元 \(x=\log r\) 将方程化为常系数方程，因而
\[
 a_n(r)=\alpha_n r^{|n|}+\beta_n r^{-|n|}.
\]
正性给出关键估计
\[
 |a_n(r)|\le\frac1{2\pi}\int_0^{2\pi}u(re^{\mathrm it})\,\mathrm dt
 =a_0(r)=O(1+|\log r|).
\]
两端乘以 \(r^{|n|}\)，再令 \(r\downarrow0\)，便知 \(\beta_n=0\)。负幂的增长比对数快，因而不能在正性所给的圆周积分界下保留。

固定 \(0<S<R\)，令 \(h_S\) 为圆周数据 \(u(Se^{\mathrm it})-B\log S\) 的 Poisson 延拓。零频率系数为 \(A\)，非零系数为 \(\alpha_n S^{|n|}\)。由\cref{b3:prop:16-poisson-frequencies}，在任意 \(0<r<S\) 上，\(h_S(re^{\mathrm it})\) 的系数恰为 \(A\) 及 \(\alpha_n r^{|n|}\)。这些也正是 \(u(re^{\mathrm it})-B\log r\) 的全部系数。

两个连续圆周函数的所有 Fourier 系数相同，就必相等。为提醒这一唯一性的依据，第16章圆环构造已证明连续周期函数的 Fej\'er 平均一致收敛；如果差函数的全部系数为零，它的每个 Fej\'er 平均都为零，因此差函数为零。于是 \(u-B\log|z|=h_S\) 于 \(0<|z|<S\)。不同 \(S\) 的调和延拓在交叠的去心圆盘一致，连续性使原点值也相同，故拼成整个 \(B(0,R)\) 上的调和函数 \(h\)。取 \(c=-B\ge0\) 得到表示。零频率确定 \(B\)，其余函数由差确定，所以表示唯一。
\end{proof}

这个结论也称正调和函数的\term[bocher-gulixing-fenjie]{Bôcher 孤立奇性分解}[Bôcher decomposition of an isolated singularity]，这里证明的是二维版本。若允许 \(u\) 改变符号，\(\operatorname{Re}(1/z)\) 就有负幂奇性；它的圆周平均为零，绝对值积分却很大，证明中的 \(|a_n|\le a_0\) 不再成立。

\Needspace{6\baselineskip}
\begin{proposition}[有限穿孔平面无正 Green 核]\label{b3:prop:17-finite-punctures-parabolic}
设 \(A\subset\mathbb C\) 为有限集，允许为空。则在 \(D=\mathbb C\setminus A\) 上，对任何 \(w\in D\) 都不存在正基本核。
\end{proposition}
\begin{proof}
空集情形已由圆周平均证明。设 \(A=\{a_1,\ldots,a_m\}\)，假设 \(q\in\mathcal Q_D(w)\)。在每个 \(a_j\) 的一个不含其余缺口及 \(w\) 的小圆盘内，上一引理给
\[
 q(z)=-c_j\log|z-a_j|+h_j(z),\quad c_j\ge0,
\]
其中 \(h_j\) 跨过 \(a_j\) 调和。因此
\[
 H(z)=q(z)+\log|z-w|+\sum_{j=1}^m c_j\log|z-a_j|
\]
延拓为整平面调和函数：在每个被删点，恰由相应局部分解消去奇性；在 \(w\) 则由正基本核的定义消去奇性。

取 \(r\) 大于 \(|w|,|a_1|,\ldots,|a_m|\)。对任意 \(|a|<r\)，有
\[
 \frac1{2\pi}\int_0^{2\pi}\log|re^{\mathrm it}-a|\,\mathrm dt=\log r.
\]
这可由 \(\log|re^{\mathrm it}-a|=\log r+\operatorname{Re}\log(1-ae^{-\mathrm it}/r)\) 的一致收敛幂级数直接积分，所有非零频率平均为零。于是
\[
 \frac1{2\pi}\int_0^{2\pi}q(re^{\mathrm it})\,\mathrm dt
 =H(0)-\left(1+\sum_{j=1}^m c_j\right)\log r.
\]
右端最终为负，左端始终为正，矛盾。
\end{proof}

每个缺口允许新增的系数均为非负，它们只能使大圆平均中的负对数斜率更大。因而不能靠有限个孤立缺口补偿无穷远的正性失败。这也是论证中必须分析缺口奇性，而不能只说“孤立点没有影响”的原因。

\subsection{\texorpdfstring{\(\mathbb C\)}{C} 的抛物性}
\label{b3:subsec:170203}

从势论角度，在任何极点都没有正基本核的区域称为\term[paowuxing-yu]{抛物型域}[parabolic domain]。这里的名称针对 Green 核的不存在性，与时间演化方程的抛物分类无关。全平面及有限穿孔平面都是上述直接证明的例子。下一节会说明，正基本核若在一个极点存在，就在所有极点存在；因此平面区域不会出现“某些极点属于 Green 情形、另一些属于抛物情形”的混合情况。

\ProseParagraphBreak
圆盘扩大的例子已经能看出发散机制。以原点为极点的经典核为 \(g_{B(0,R)}(z,0)=\log(R/|z|)\)，这个函数直接满足上一节的全部定义条件。固定 \(z\ne0\)，随 \(R\to\infty\) 核值趋于正无穷。每个有界圆盘都可容纳有限点源响应，但把零边界不断推向无穷远后，在固定观察点得不到有限极限。

\ProseParagraphBreak
维数在这里有实际作用。在三维，正基本解与 \(|x-w|^{-1}\) 成比例，它在整个空间保持正且在无穷远趋零；二维的基本解则是对数。因而“把边界推到无穷远”是否仍保留有限正核，与基本解的尺度有关，不能把平面的圆周平均论证原样移到高维。

\subsection{覆盖类型与 Green 存在性的区别}
\label{b3:subsec:170204}

Riemann 曲面论通常把万有覆盖为圆盘的曲面称为双曲型；有些势论文献也用“双曲型”表示存在正 Green 函数。两种说法描述不同条件，本章统一用“Green 域”指核的存在性。

\ProseParagraphBreak
去心圆盘是有界域，所以是 Green 域；它的万有覆盖也是圆盘。另一方面，\cref{b3:prop:12-planar-hyperbolic}说明，补集至少含两个点的平面域具有圆盘万有覆盖，故 \(\mathbb C\setminus\{0,1\}\) 属于这种覆盖类型；该结论沿用第12章已标明的一般统一化输入。但\cref{b3:prop:17-finite-punctures-parabolic}表明它没有正 Green 核。这个例子同时具有圆盘覆盖与势论抛物性。已知万有覆盖的共形类型，尚不足以判断点源在原域内是否产生有限正响应。

\ProseParagraphBreak
单连通情形则可以直接转移。若 \(D\subsetneq\mathbb C\) 为单连通域，Riemann 映照定理给出共形双射 \(F:D\to\mathbb D\)。固定 \(w\in D\)，取圆盘的经典 Green 核，并令
\[
 q(z)=g_{\mathbb D}(F(z),F(w)).
\]
它在极点外正且调和。因 \(F'(w)\ne0\)，在 \(w\) 附近有
\[
 F(z)-F(w)=(z-w)A(z),\quad A(w)=F'(w)\ne0,
\]
其中 \(A\) 全纯且局部无零。因此复合核加上 \(\log|z-w|\) 后，只新增调和项 \(-\log|A(z)|\)，仍具有所需奇性。它是 \(\mathcal Q_D(w)\) 的成员，故任意真单连通平面域都是 Green 域。这个存在性判断只用内部共形双射，不要求 \(F\) 连续延拓到通常拓扑边界。

\ProseParagraphBreak
若只有覆盖映射而没有双射，原域的一个极点会有多个原像。把原域的核拉回覆盖空间，便在每一个原像处产生对数奇性，不能把这个拉回误认为覆盖空间上只有一个极点的 Green 核。这正是多连通情形比共形双射多出的困难。

\subsection{与万有覆盖的联系}
\label{b3:subsec:170205}

设 \(\pi:\mathbb D\to D\) 为万有覆盖，固定 \(w\in D\) 及一个原像 \(\widetilde w\)。第11章的覆盖理论说明，纤维由覆盖变换群 \(\Gamma\) 的轨道 \(\{\gamma\widetilde w:\gamma\in\Gamma\}\) 给出；该集合离散且可数。于是一个自然的候选表达是
\[
 \sum_{\gamma\in\Gamma}
 g_{\mathbb D}(\widetilde z,\gamma\widetilde w).
\]
每项都非负；有限部分和只处理有限多个点源，而完整纤维要求控制无穷项。即使每个单项都有限，和仍可能发散。离散性只保证内部紧集包含有限个极点，并不控制趋近单位圆周的无穷多个极点的总贡献。

\ProseParagraphBreak
要用这一表达构造原域核，至少还须证明：部分和在极点轨道外局部一致收敛；在每个原像附近只留下一个系数为一的对数奇性；和对覆盖变换不变，从而能够下降。要进一步识别为最小核，还需最小性比较。这些问题各有作用，不能用“按群求和”一句话取代。本章的一般构造采用域内穷竭，因而不以该级数公式为前提。

\ProseParagraphBreak
从这里也可理解为何前述两种分类会分开。圆盘覆盖给出每个单独源点的正核，却没有自动给出整个轨道的求和控制；势论中的有限性正是额外要求。下一节把这种要求改写成更便于检验的条件：在一个固定非极点位置，穷竭核的数值有界。Harnack 链随后会把这个一点条件传播到所有内部紧集。
